\documentclass{aastex701}

\begin{document}

\title{A Giant Transiting Planet Candidate on a 6.26 Day Orbit around the A1V Star TIC 264626874}

\author[0000-0000-0000-0000]{Carl Drinkwater}
\affiliation{Independent researcher, Shropshire, UK}
\email{carl@29degrees.co.uk}

\begin{abstract}
I report a transit signal on TIC 264626874 (TYC 0113-01567-1, $T=10.07$), an A1V star of 2.9 solar radii.
I found it in a box least squares search of the Quick-Look Pipeline light curves of TESS Sector 98, and recovered it in
the TESS-SPOC light curve of Sector 32. A joint fit of both sectors gives a period of 6.2599677 d, a depth of 3136 ppm
and a radius of $17.7 \pm 0.9\,R_\oplus$. Difference images place the transit on the target. The light curves show no
secondary eclipse and no significant odd-even depth difference. TRICERATOPS gives a false positive probability of
$0.17 \pm 0.13$, with the planet on the target the most likely scenario. The radius does not separate a giant planet from
a brown dwarf or a low-mass star, so radial velocity or Doppler tomography follow-up is needed.
\end{abstract}

\section{Detection and vetting}

I searched the Quick-Look Pipeline \citep[QLP;][]{huang2020a,huang2020b,kunimoto2021} light curves of TESS
\citep{ricker2015} Sector 98 for stars brighter than $T=11.5$. I removed trends with a biweight filter of 0.5 d
\citep{hippke2019} and searched 0.5 to 14 d with box least squares \citep{kovacs2002}. TIC 264626874 \citep{stassun2019}
shows a signal at 6.2596 d with a multiple event statistic of 38.9. It passes the flux-level false alarm and false
positive tests of LEO-vetter \citep{kunimoto2025}, and matches no TOI \citep{guerrero2021}, CTOI, SPOC threshold
crossing event or TESS eclipsing binary \citep{prsa2022} in the lists of 2026 October 16.

The TESS-SPOC light curve of Sector 32 \citep{caldwell2020} shows the same transit at a signal-to-noise ratio of 45.9.
TESS observed the star in no other released sector. A joint box least squares search of both sectors finds the
signal at 6.25999 d with a signal-to-noise ratio of 62.5, against 16.7 for the highest of the null trials.

\section{Transit fit}

I fitted a transit model \citep{kreidberg2015} to both sectors with emcee \citep{foremanmackey2013}, with the cadence
of each sector. Figure~\ref{fig:transit} shows the fit. The fit gives:
period $6.2599677 \pm 0.0000039$ d;
mid-transit time $2461011.98304 \pm 0.00095$ BJD$_\mathrm{TDB}$;
depth $3136 \pm 45$ ppm;
duration $5.483 \pm 0.070$ h;
$R_p/R_\star = 0.0560 \pm 0.0006$;
impact parameter $0.745 \pm 0.033$;
$a/R_\star = 6.58 \pm 0.35$.
With the TIC stellar radius of $2.90 \pm 0.14\,R_\odot$, the companion radius is $17.7 \pm 0.9\,R_\oplus$
($1.58\,R_\mathrm{J}$). The stellar density from the fit is 0.97 of the TIC value.

The depths of the pipeline light curves differ: $2770 \pm 57$ ppm in Sector 98 and $3385 \pm 45$ ppm in Sector 32.
Aperture photometry with the same apertures on the TESScut cutouts \citep{brasseur2019} of both sectors gives 2694 to
3096 ppm, consistent between sectors ($p = 0.70$). I therefore attribute the difference to the background and
dilution corrections of the two pipelines.

\section{False positive tests}

The odd and even transits differ by 1.3$\sigma$. The flux at phase 0.5 is $-18 \pm 34$ ppm, against a maximum of 157 ppm
for the thermal and reflected light of a planet. The difference images \citep{bryson2013} of Sectors 32 and 98 place
the transit on the target, at signal-to-noise ratios of 10.9 and 7.7 (13.4 jointly). Two Gaia DR3 stars within one
arcminute could cause the dip only with eclipses of 24\% and 27\%, and the difference images exclude both as the
source.

Gaia DR3 \citep{gaia2023} gives a RUWE of 1.19, a non\_single\_star flag of 0, and no radial velocity variability
in 31 epochs ($p = 0.51$; \citealt{katz2023}). The epoch scatter of 11.1 km\,s$^{-1}$, from the line broadening of
185 km\,s$^{-1}$, limits the companion mass only to below 0.26 $M_\odot$. A LAMOST spectrum \citep{cui2012} classifies the
star as A1V, consistent with the TIC temperature of 9600 K.

I fitted the ellipsoidal, reflection and beaming terms \citep{faigler2011} to the light curves before detrending.
The ellipsoidal amplitude is $163 \pm 107$ ppm, where the error comes from the same fit at 200 wrong periods. A
companion of 0.1 $M_\odot$ would give about 160 ppm and a Jupiter-mass planet about 1.5 ppm, so this measurement
excludes neither.

TRICERATOPS \citep{giacalone2021} gives a false positive probability of $0.17 \pm 0.13$ and a nearby false positive
probability of 0 over five runs, with a transiting planet on the target as the most likely scenario.

\section{Conclusion}

TIC 264626874 hosts a transiting companion of $17.7 \pm 0.9\,R_\oplus$ on a 6.26 d orbit. The companion may be an
inflated giant planet, a brown dwarf or a low-mass star. The host is bright and rotates rapidly, so Doppler
tomography can confirm that the shadow of the transit is on the star, and radial velocities can measure the mass.
TESS observes the star again in Sector 110.

\begin{figure}[ht!]
\plotone{figure_1.pdf}
\caption{Transit of TIC 264626874 in the Sector 98 QLP (left) and Sector 32 TESS-SPOC (middle) light curves, with the
joint fit (red) and 37-minute bins (black). Right: the flux at phase 0.5, with the transit duration shaded.
\label{fig:transit}}
\end{figure}

\begin{acknowledgments}
This note uses data from the TESS mission, obtained from MAST. Claude (Anthropic) wrote the analysis code and a draft of
this note under my direction. The code and results are at https://github.com/carldr/tess-transiting-planet-search.
\end{acknowledgments}

\facilities{TESS, Gaia, LAMOST}

\software{astropy, astroquery, batman \citep{kreidberg2015}, emcee \citep{foremanmackey2013}, LEO-vetter
\citep{kunimoto2025}, lightkurve, TRICERATOPS \citep{giacalone2021}, wotan \citep{hippke2019}}

\bibliography{references}{}
\bibliographystyle{aasjournalv7}

\end{document}
